% Generated draft: requires mathematical and source review.
\documentclass[11pt]{article}
\usepackage[margin=25mm]{geometry}
\usepackage{fontspec}
\setmainfont{Latin Modern Roman}
\usepackage{amsmath,amssymb,mathtools,unicode-math}
\setmathfont{Latin Modern Math}
\usepackage{enumitem,needspace,xurl,hyperref}
\setlength{\parindent}{0pt}
\setlength{\parskip}{5pt}
\setlength{\emergencystretch}{3em}
\newcommand{\N}{\mathbb{N}}
\newcommand{\Z}{\mathbb{Z}}
\newcommand{\Q}{\mathbb{Q}}
\newcommand{\R}{\mathbb{R}}
\newcommand{\C}{\mathbb{C}}
\newcommand{\F}{\mathbb{F}}
\newcommand{\A}{\mathbb{A}}
\newcommand{\PP}{\mathbb{P}}
\newcommand{\E}{\mathbb{E}}
\newcommand{\eps}{\varepsilon}
\newcommand{\dd}{\,\mathrm{d}}
\begin{document}
\section*{UnsolvedMath ID: 30001300}
\textbf{Source problem code:} catalogue entry
\subsection{Definitions and mathematical statement}
The supplied problem asks for a higher-dimensional analogue of the curve condition \(g>1\), together with a generalization of an unspecified crude height-counting heuristic that would predict the Bombieri--Lang conjecture for all varieties of general type. The following gives a conventional precise conjectural target; its field and geometric hypotheses are explicit interpretive choices because the supplied statement does not specify them.

Let \(k\) be a number field and let \(X/k\) be a smooth, projective, geometrically integral variety of dimension \(n\geq 1\). Write \(K_X\) for its canonical line bundle. Say that \(X\) is of general type if \(K_X\) is big, meaning
\[
\limsup_{m\to\infty}\frac{\dim_k H^0(X,K_X^{\otimes m})}{m^n}>0.
\]
For a smooth projective geometrically integral curve \(C/k\), this condition is equivalent to \(g(C)>1\). Thus general type supplies the conventional higher-dimensional analogue of the stated curve condition.

Under this interpretation, the conjectural prediction to be explained is
\[
\text{for every number field }k\text{ and every such }X/k\text{ of general type},\qquad
\overline{X(k)}^{\mathrm{Zar}}\subsetneq X.
\]
Equivalently, for each such pair \((k,X)\), there exists a proper Zariski-closed subset \(Z\subsetneq X\), defined over \(k\), such that \(X(k)\subseteq Z(k)\). The subset \(Z\) may depend on \(k\) and \(X\). This predicts non-density, without asserting finiteness of \(X(k)\) in higher dimensions.

To express the counting aspect, after fixing an ample line bundle \(L\) on \(X\) and an associated absolute multiplicative Weil height \(H_L\), one may consider
\[
N_{X,k,L}(B)=\#\{P\in X(k):H_L(P)\leq B\},\qquad B\geq 1.
\]
The source question is to formulate a higher-dimensional height-counting heuristic that predicts the displayed non-density statement for every variety of general type under the intended hypotheses. The supplied record gives neither the original crude heuristic nor a proposed counting law, height normalization, or treatment of exceptional subvarieties. Consequently, it does not determine a unique precise heuristic or asymptotic assertion. In particular, a prediction of slow growth for \(N_{X,k,L}(B)\) alone is not a statement of Zariski non-density. These unresolved qualifications are part of this draft; no heuristic or proof is supplied.
\subsection{Short English statement}
General type is the conventional higher-dimensional analogue of \(g>1\) for smooth projective curves. The question asks for a height-counting heuristic predicting that rational points on every general-type variety are contained in a proper closed subset. The displayed formulation uses explicit number-field and smooth projective hypotheses; the source does not specify these hypotheses or the heuristic to be generalized.
\subsection{Sources}
\begin{itemize}
\item Local UnsolvedMath registry, supplied complete record: permanent ID \(30001300\), exact source problem code \texttt{catalogue entry}, title ``Height Heuristics for Bombieri--Lang in Higher Dimensions,'' in the Oberwolfach Reports Open Problems collection. The record lists proposed year \(2009\) and supplies no named proposer.
\item Supplied source citation: \emph{Explicit Methods in Number Theory (2010)}, Oberwolfach Reports, DOI \texttt{10.4171/owr/2009/33}, pp.~1843--1920. \url{https://doi.org/10.4171/owr/2009/33}. This citation is reproduced as supplied; its year information has not been reconciled or independently checked.
\item Supplied literature source labelled ``Poonen Oberwolfach report'': \url{https://ems.press/content/serial-article-files/46233?nt=1}. This label does not establish attribution of the particular problem to Poonen.
\item Supplied literature source labelled ``Bombieri--Lang survey'': \url{https://arxiv.org/abs/2012.15765}.
\item Literature assessment dated exactly \texttt{2026-08-21}: the supplied curation labels the problem open and describes higher-dimensional Bombieri--Lang non-density as a major open conjecture with substantial special cases proved. This is supplied dated curation derived from the OWR material, not an independently verified literature assessment. The record requests independent verification of time-sensitive status and cited claims.
\item The supplied statement-status label is \texttt{corrected\_verified}, with an accompanying claim of meaning-preserving semantic, typographic, and LaTeX cleanup. This records the supplied editorial assessment, not mathematical verification of this draft.
\end{itemize}
\end{document}
